\section{Kịch bản vận hành đầu cuối}

Kịch bản đánh giá xuyên suốt liên kết kết quả tối ưu mô hình với nền tảng:

\begin{enumerate}
  \item Đẩy engine của mô hình tối ưu lên kho model; phân giải thành URL đã ghim
        commit.
  \item Tạo cluster, đăng ký node; khởi động Edge Agent trên thiết bị và xác nhận
        node chuyển ONLINE trên Dashboard.
  \item Tạo deployment AI Runtime với \texttt{model\_uri} vừa tạo, nguồn là camera
        hoặc video công trường.
  \item Theo dõi FPS, latency, số detection, confidence trên Dashboard qua
        WebSocket.
  \item Dừng container runtime thủ công trên thiết bị và quan sát vòng reconcile
        tạo lại container.
  \item Hỏi AI Assistant về trạng thái node và telemetry gần nhất; yêu cầu
        Assistant khởi động lại runtime và xác nhận.
  \item Cập nhật deployment sang model khác, sau đó cập nhật ngược lại model cũ
        để khôi phục.
  \item Tắt Edge Agent và xác nhận Cloud đánh dấu node OFFLINE.
\end{enumerate}

\begin{table}[H]
\centering
\caption{Kết quả kịch bản đầu cuối}
\label{tab:e2e-scenario}
\begin{tabularx}{\textwidth}{L{5.4cm}C{2.6cm}Y}
\toprule
\textbf{Tiêu chí} & \textbf{Kiểm thử tự động} & \textbf{Trên thiết bị}\\
\midrule
Node đăng ký và báo ONLINE & Đạt & \cbs{}\\
Triển khai runtime và báo SUCCEEDED & Đạt & \cbs{}\\
Tự tạo lại container bị dừng & Đạt & \cbs{}\\
Telemetry hiển thị realtime & Đạt (WebSocket) & \cbs{}\\
Assistant tra cứu và điều khiển có xác nhận & Đạt (đọc) & \cbs{}\\
Cập nhật và khôi phục model & Đạt (update) & \cbs{}\\
Phát hiện node OFFLINE & Theo thiết kế & \cbs{}\\
\bottomrule
\end{tabularx}
\end{table}

Một lần chạy được xem là thành công khi trạng thái trên Cloud khớp trạng thái
thực tế trên thiết bị ở mọi bước, không cần thao tác SSH vào thiết bị ngoài bước
chủ động gây lỗi.
